\section{Introduccion}
\label{sec:introduction}

Las alteraciones neu\-ro\-ló\-gi\-cas abarcan un conjunto heterogéneo de trastornos que impactan tanto el sistema nervioso central como el periférico, comprometiendo funciones críticas como la motricidad, memoria, cognición y percepción. Estas patologías constituyen una de las causas predominantes de discapacidad e incapacidad a nivel mundial \cite{Ding2022}. Durante las últimas décadas, su prevalencia ha experimentado un incremento significativo asociado al envejecimiento poblacional y a múltiples factores de riesgo concurrentes, generando una demanda creciente de instrumentos diagnósticos y terapéuticos mejorados \cite{Ding2022}. La heterogeneidad clínica de estas afecciones —que incluye procesos neurodegenerativos, cerebrovasculares y funcionales— complica considerablemente el diagnóstico debido a la complejidad y sobreposición de sus presentaciones clínicas \cite{Mavroudis2024} \cite{Finkelstein2025}.

Entre las patologías oncológicas del sistema nervioso central, los tumores cerebrales representan un grupo diverso de neoplasias con origen primario encefálico o secundario por metastatización, exhibiendo grados de malignidad e historias naturales muy variables \cite{Pichaivel2022}. El glioblastoma destaca como la neoplasia maligna primaria intracraneal más prevalente en la población adulta, caracterizándose por su conducta infiltrativa, alta tendencia a recidivar y pronóstico desfavorable aun con tratamientos contemporáneos \cite{Schaff2023}. La sintomatología clínica presenta variabilidad en función de la localización topográfica, magnitud y velocidad de expansión tumoral, frecuentemente manifestándose mediante cefaleas refractarias, crisis convulsivas, disfunción cognitiva y déficits neurológicos focales que deterioran progresivamente la calidad de vida \cite{Barkade2023} \cite{Essianda2023}. La complejidad inherente a la biología molecular y fisiopatología de estas lesiones constituye un obstáculo significativo para el diagnóstico preciso y la selección óptima de estrategias terapéuticas \cite{Essianda2023}.

La resonancia magnética (MRI) permanece como la modalidad de neuroimagen de elección para la detección y caracterización de masas intracraneales, gracias a su excelente resolución de contraste en estructuras de partes blandas y su capacidad para determinar con precisión la localización espacial, extensión y características morfológicas de lesiones \cite{Martucci2023}. En paralelo, el desarrollo de metodologías radiómicas permite extraer atributos cuantitativos de las imágenes para mejorar la clasificación de neoplasias, estimación pronóstica y orientación clínica mediante la integración de técnicas de inteligencia artificial \cite{Wang2025}. Las aproximaciones modernas en neuroimagen incorporan información anatómica, dinámica y metabólica de forma integrada, facilitando una comprensión más completa de la naturaleza tumoral y optimizando la planificación de intervenciones terapéuticas \cite{Nikam2022}. En esta perspectiva, la segmentación automática de imágenes de resonancia magnética se ha consolidado como un componente fundamental para delimitar territorios tumorales y sustentar sistemas de asistencia diagnóstica basados en aprendizaje automático y redes neuronales profundas \cite{Mohammed2023}.

La tecnología de resonancia magnética por imágenes (MRI) constituye un método diagnóstico no invasivo que aprovecha campos magnéticos intensos y radiaciones de radiofrecuencia para producir representaciones de alta definición de órganos y tejidos, sin exposición a radiación ionizante \cite{AhmedAbdullahAbdulazizALRabah2024}\cite{Taha2024}. La superior resolución espacial combinada con el contraste diferencial entre tejidos blandos la posiciona como una herramienta indispensable en la evaluación de padecimientos neurológicos, particularmente en la visualización y caracterización de masas encefálicas \cite{AhmedAbdullahAbdulazizALRabah2024}\cite{Hussain2022}. Los avances recientes en protocolos de adquisición y algoritmos de procesamiento han mejorado substancialmente la calidad diagnóstica y ampliado el espectro de aplicaciones clínicas \cite{Alsaedi2024}. Por estas razones, la MRI constituye la base fundamental de numerosos sistemas computacionales de inteligencia artificial, suministrando imágenes de elevada resolución que facilitan la detección y clasificación automatizada de irregularidades del sistema nervioso central \cite{Alsaedi2024}\cite{Hussain2022}\cite{Taha2024}.

Un aspecto limitante del enfoque tradicional radica en que la inspección visual de imágenes resulta sumamente laboriosa y vulnerable a la fatiga atencional de profesionales en radiología \cite{Abdusalomov2023}. Además, la identificación manual es propensa a imprecisiones derivadas de la variabilidad en la interpretación humana. La necesidad de analizar decenas de cortes transversales para emitir un juicio diagnóstico confiable genera un volumen de trabajo que puede retrasar el inicio de intervenciones médicas urgentes.

La identificación de masas encefálicas presenta complejidad inherente a su variabilidad morfológica. Las neoplasias cerebrales típicamente exhiben configuraciones asimétricas, límites difusos e intensidades de señal similares al tejido nervioso sano. La presencia concomitante de artefactos de imagen o degradación de la calidad por error técnico dificulta adicionalmente la visibilidad de hallazgos patológicos \cite{Soomro2023}. Las metodologías convencionales de análisis frecuentemente fallan ante estas variaciones en iluminación y contraste. Estas limitaciones fundamentan la necesidad clínica de implementar sistemas de procesamiento automatizado con elevada precisión diagnóstica.

La inteligencia artificial emerge como una alternativa promisoria ante estos desafíos clínicos. Los sistemas computacionales poseen la capacidad de procesar volúmenes masivos de imágenes en tiempos muy breves. Estas herramientas funcionan como soporte de segunda opinión para especialistas en radiología en contextos hospitalarios. El empleo de algoritmos sofisticados minimiza significativamente la inconsistencia inter e intraobservador en valoraciones diagnósticas. Esta tecnología está revolucionando los enfoques actuales para patologías oncológicas \cite{Anantharajan2024}.

El aprendizaje profundo ha demostrado superioridad en tareas de reconocimiento visual médico. Las Redes Neuronales Convolucionales (CNN) representan el paradigma establecido para estas labores perceptivas complejas. Estos modelos matemáticos extraen de forma completamente automatizada los descriptores anatómicos presentes en imágenes de resonancia, sin necesidad de ingeniería manual de atributos \cite{Khan2022}. Las CNN eliminan la dependencia de identificar manualmente áreas de interés clínico en las imágenes. Su potencial analítico potencia la detección de anomalías estructurales encefálicas.

Numerosas investigaciones sustentan la alta performance de estas arquitecturas en escenarios clínicos experimentales \cite{Mahmud2023}. Los modelos contemporáneos logran diferenciar de manera efectiva entre diversas categorías de tumores cerebrales. La aplicación de estas metodologías acelera significativamente los tiempos de respuesta en procesos decisorios clínicos \cite{Dorfner2025}. No obstante, la evaluación sistemática y comparativa de diferentes arquitecturas bajo condiciones equivalentes permanece como una necesidad de investigación. Es fundamental contrastar el desempeño de estos algoritmos empleando bases de datos públicas y recientes.

El objetivo central de esta investigación es realizar una comparación estructurada de múltiples modelos de aprendizaje profundo para la identificación de tumores en estudios de resonancia magnética cerebral. El estudio analiza el rendimiento de distintas arquitecturas neuronales convolucionales utilizando un dataset obtenido de la plataforma Kaggle. La presentación del artículo sigue una estructura progresiva. En la sección siguiente se exponen los trabajos previos identificados mediante herramientas de inteligencia artificial. Las secciones posteriores detallan de forma sistemática los materiales y procedimientos experimentales, los hallazgos de cada arquitectura, el análisis integrado, la discusión interpretativa y las conclusiones derivadas.
